%% Card-local preamble: the source-import-staged retained-preamble.tex of this unit is reproduced verbatim below (author macro and theorem block plus the publisher-class abbreviations it needs).
\usepackage{amsmath,amssymb,amsthm,mathtools}
\usepackage{xcolor}
\usepackage{url}
\usepackage[hidelinks]{hyperref}
\usepackage{xfrac}

\numberwithin{equation}{section}

\newtheorem{Theorem}{Theorem}[section]
\newtheorem*{Theorem*}{Theorem}
\newtheorem{Corollary}[Theorem]{Corollary}
\newtheorem{Lemma}[Theorem]{Lemma}
\newtheorem{prop}[Theorem]{Proposition}
 { \theoremstyle{definition}
\newtheorem{Definition}[Theorem]{Definition}
\newtheorem{Note}[Theorem]{Note}
\newtheorem{Example}[Theorem]{Example}
\newtheorem{Remark}[Theorem]{Remark} }
\newtheorem{rhp}[Theorem]{Riemann--Hilbert Problem}

\usepackage[mathscr]{euscript}
\usepackage{commath, enumitem, tikz}
\newcommand*\up{\ensuremath{\vcenter{\hbox{\includegraphics[width=.6em]{Figures/squareur.pdf}}}}}

\newcommand*\lw{\ensuremath{\vcenter{\hbox{\includegraphics[width=.6em]{Figures/squarell.pdf}}}}}

\newcommand{\ee}{\mathrm{e}}
\newcommand{\dd}{\mathrm{d}}
\newcommand{\ii}{\mathrm{i}}
\newcommand{\para}[1]{\breve{#1}}

\newcommand{\R}{{\mathbb{R}}}
\newcommand{\N}{{\mathbb{N}}}
 \newcommand{\C}{{\mathbb{C}}}
\newcommand{\Q}{{\mathbb{Q}}}
\newcommand{\Z}{{\mathbb{Z}}}
\newcommand\lb{\left(}
\newcommand\rb{\right)}

\newcommand{\re}{\operatorname{Re}}
\newcommand{\im}{\operatorname{Im}}
\newcommand{\Tr}{\operatorname{Tr}}
\newcommand{\rhy}{{\rm RHP}-${\mathbf Y}$}

\newcommand{\ic}{{\mathrm i}}

\newcommand{\supp}{\operatorname{supp}}
\renewcommand{\deg}{\operatorname{deg}}
\renewcommand{\arg}{\operatorname{Arg}}

\newcommand{\sgn}{\operatorname{sgn}}
\newcommand{\qandq}{\qquad \text{and} \qquad}
%\newcommand{\qasq}{\quad \text{as} \quad }
%\newcommand{\qforq}{\quad \text{for} \quad }
%\newcommand{\br}{\textcolor{brown} }

